(1) 沿左向右定义正热流 \(J=-\kappa(T)T_x\)。左热右冷时左A层、右B层满足
\[
JL=\int_{T_i}^{T_h}aT\,dT=\tfrac a2(T_h^2-T_i^2),
\quad JL=\int_{T_c}^{T_i}b\,dT=b(T_i-T_c).
\]
两式相等构成严格单调界温方程：前者随 \(T_i\) 增大而减少，后者随 \(T_i\) 增大而增加，所以 \(T_c<T_i<T_h\) 内唯一交点，不可用平均温度替代整个变导热积分。

(2) 右热左冷时热从右往左，穿过B层先由 \(T_h\) 降至新界温 \(T_i'\)，再穿过A降至 \(T_c\)，故热流大小满足
\[
|J'|L=b(T_h-T_i')=\tfrac a2(T_i'^2-T_c^2).
\]
与第一问界面温度方程不同，这是非线性材料在不同温区的组合所致。若把两个材料连同热源也同时交换，仅仅重新标记坐标则当然不会产生整流；本题比较的是同一器件交换冷热端。

(3) 令单位 \(a=b=L=1\)，正向方程 \((4-T_i^2)/2=T_i-1\)，得 \(T_i=-1+\sqrt7\)，正向热流 \(\boxed{J_{\rightarrow}=\sqrt7-2}\)。
反向 \(2-T_i'=(T_i'^2-1)/2\)，得 \(T_i'=-1+\sqrt6\)，反向热流 \(\boxed{J_{\leftarrow}=3-\sqrt6}\)。数值分别约为 \(0.64575\) 与 \(0.55051\)，整流比约 \(1.173>1\)。两界温都位于1与2之间，确保两层热流符号一致而没有违反Fourier定律。

(4) 每一层的熵产生率密度为 \(\sigma=\kappa(T)(T_x)^2/T^2=J^2/(\kappa T^2)\ge0\)，积分并用 \(dx=-\kappa dT/J\) 得单层 \(\Sigma_{\rm layer}=|J|(1/T_{\rm cold,layer}-1/T_{\rm hot,layer})\)。两层界温项严格抵消，总熵生产
\[
\boxed{\Sigma=|J|\left(\frac1{T_c}-\frac1{T_h}\right)\ge0}
\]
与冷热热库总熵变相等。若 \(a\to0\)，A近绝热，两方向热流都趋零而非产生无限整流；\(T_h-T_c\to0\) 时在同一个平均温度线性化，两个方向的首阶导热相同，整流从较高阶温差才显现。
